\scanpage{78}
% loc: 1–3, boxed upper right
\fbox{\parbox{.92\linewidth}{Description \deleted{du} \inserted{de l'anneau} $\mathcal A\simeq M_2(k)$ ($k$ corps) en termes \unsure{de|du} $G=\mathcal A^\times/k^\times$ (tentative)}}

% loc: 1–5, left
\[
\begin{aligned}
1*1&=1\\
(x\deleted{\unsure{s|a}}*y)(y*x)&=1\\
xs*y&=x*y\quad\text{si } s,y\text{ commutent}\\
[\,x*yt&=x*y\quad\text{si }t,x\quad\text{---}\,]\\
[\,x*y&=1\quad\text{si } x,y\text{ commutent}\,]
\end{aligned}
\]

% loc: 3–5, right
\[
\begin{aligned}
[\,x*y&=\bar{x}\,\bar{y}\,\bar{x}^{-1}\bar{y}^{-1}\\
\delta(\dot x)&=\deleted{\unsure{\mathrm{discriminant}|\mathrm{discrim}}}\frac{b^2-4c}{c}\,]\\
c&=\det\dot x,\ b=\operatorname{Tr}x
\end{aligned}
\]

% loc: 6–7
\[
(x*y)^2-P(\delta x,\delta y,\delta(xy))\,x*y+1=0\qquad P\in\mathbb Z[X,Y,\unsure{Z|\mathbb Z}]
\]
\margin{bien déterminé}

\[
[\delta(z),x*y]=[\delta(z),\delta(z')]=0
\]

% loc: 8–11, framed group and annotation below
\[
\boxed{\left\{\begin{aligned}
\delta(xy)&=\delta(yx)\qquad\text{i.e. }\delta(g(x))=\delta(x)\qquad(x,y,g\in G)\\
\delta(x)&=\delta(x^{-1})\\
\delta(xy)+\underbrace{\delta(xy^{-1})}_{\shortparallel\,\delta(x^{-1}y)}
 &=\deleted{\unsure{\delta(x)\delta(y)}}(\delta(x)+2)(\delta(y)+2)-4
 =\delta(x)\delta(y)+2(\delta(x)+\delta(y))
\end{aligned}\right.}
\]
\margin{si $x,y$ commutent seulement ??? À vérifier !}

% loc: 12–13
\[
[\,\boxed{\delta(1)=0}
\]
\[
\delta(x^n)=P_n(\delta(x)),\qquad\delta(\underbrace{x^ny^m}_{\text{ou d'autres monômes non commutatifs\ldots}})=P_{n,m}(\delta(x),\delta(y),\delta(xy))\quad\text{etc.}
\]

% loc: 14–18
$x_1\ x_2\ x_3\ x_4,\in G$, il doit y avoir une relation pol. à coefficients dans $\mathbb Z$ entre
\[
\delta(x_i)\quad(1\leq i\leq4),\quad\deleted{\unsure{\delta}}\quad\delta(x_ix_j)\quad(1\leq i<j\leq4)
\]
(dix quantités, \deleted{\unsure{pour en un ``espace''|pour sur un ``espace''|pour un ``exprimant''}} dépendant de ``neuf paramètres''\ldots\ ]

% loc: 19–20, boxed
\[
\boxed{\delta(xyz)=\underset{\text{pol. univ. à coeff. entiers en six variables}}{Q}(\delta(x),\delta(y),\delta(z),\delta(xy),\delta(yz),\delta(zx))\qquad ?}
\]
\margin{à déterminer}

% loc: 21–24
\deleted{Ceci donne $\delta(x^n)$ comme pol.}
\[
\begin{aligned}
\delta(x^2)\deleted{\unsure{4|+}}&=(\delta(x)+2)^2-4=\delta(x)^2+4\delta(x)\\
\delta(x^n)&=P_n(\delta(x))\qquad n\in\mathbb Z\qquad P_n=P_{-n}\\
\delta(x^ny^m)&=P_{\unsure{n,m|m,m}}(\delta(x),\delta(y),\delta(xy))
\end{aligned}
\]
\margin{[cas particulier de $\delta(xy)+\delta(xy^{-1})=\ldots$]}

% loc: 25–28
\[
\delta\bigl(\deleted{\text{monôme en}}\,\inserted{M(x_1,\ldots,x_n)}\bigr)=P_M\bigl((\delta(x_i)),(\delta(x_ix_j))\bigr)
\]
\margin{monôme non commutatif à exposants entiers $\in\mathbb Z$}
\margin{pol. à coeff. dans $\mathbb Z$ (pas uniquement déterminé si $n\geq4$)}
% archive sheet number: 77

